\begin{table}[htbp]
\centering
\caption{Covariate Balance at the Cutoff}
\label{tab:balance}
\begin{tabular}{lccccr}
\hline\hline
Covariate & RDD Estimate & SE & $p$-value & Bandwidth & Eff. $N$ \\
\hline
n Authors & 0.2375 & (0.3436) & 0.490 & 54 & 676 \\
n Categories & -0.1903 & (0.2366) & 0.421 & 33 & 426 \\
Abstract Length & -47.9551 & (70.9040) & 0.499 & 40 & 511 \\
Cs AI & -0.1134 & (0.0900) & 0.208 & 33 & 429 \\
Cs CL & 0.0533 & (0.1064) & 0.616 & 32 & 413 \\
Cs LG & -0.2462 & (0.1088) & 0.024 & 37 & 479 \\
Stat ML & -0.0855 & (0.0985) & 0.385 & 40 & 511 \\
Cs CV & -0.0028 & (0.0664) & 0.966 & 30 & 389 \\
\hline\hline
\end{tabular}
\begin{minipage}{0.95\textwidth}
\footnotesize
\textit{Notes:} Each row reports a separate RDD estimate of the discontinuity in the covariate at the 14:00 ET cutoff, using \texttt{rdrobust} with MSE-optimal bandwidth selection. Balance tests use the full arXiv metadata sample (unmatched, $N \approx 4{,}000$ within $\pm$120 minutes) rather than the citation-matched subsample, because pre-determined covariates (number of authors, categories, abstract length) are observable for all submissions regardless of citation matching. Bandwidth and effective sample size (Eff. $N$) vary across covariates because each uses its own MSE-optimal bandwidth; the larger effective samples relative to the matched outcome regressions reflect the larger pool available for balance testing.
\end{minipage}
\end{table}
